%==========================================================%
% global_config.tex — informations de couverture            %
%                                                           %
% ⚠ LES CHAMPS MARQUÉS « À COMPLÉTER » DOIVENT ÊTRE REMPLIS  %
%   AVANT LE DÉPÔT : nom de l'étudiant, encadrants.          %
%==========================================================%

%============= Types de colonnes du template ==============%
\newcolumntype{L}{>{\raggedright\arraybackslash}}
\newcolumntype{R}{>{\raggedleft\arraybackslash}}
\newcolumntype{C}{>{\centering\arraybackslash}}
%==========================================================%

%================== Couverture ============================%

\title{Conception et développement d'un moteur agentique de coaching de vente et d'optimisation dynamique des stocks pour le Retail en temps réel}

\author{Prénom NOM}                    % ← À COMPLÉTER

%%% Si binôme, décommenter les deux lignes suivantes
%\setboolean{isBinomal}{true}
%\secondAuthor{Prénom NOM}

\diplomaName{Diplôme National d'Ingénieur en Sciences Appliquées et Technologiques}
\speciality{Science des Données et Intelligence Artificielle}

%% Encadrant professionnel
\proFramerName{Monsieur Prénom NOM}    % ← À COMPLÉTER
\proFramerSpeciality{Encadrant professionnel, Ooredoo Tunisie}

%% Encadrant académique
\academicFramerName{Monsieur Prénom NOM} % ← À COMPLÉTER
\academicFramerSpeciality{Encadrant académique, TEK-UP}

%% Entreprise d'accueil
\companyName{Ooredoo Tunisie}

%% Année universitaire
\collegeYear{2025 - 2026}

%%%%%% Signatures %%%%%%

\proSignSentence{J'autorise l'étudiant à faire le dépôt de son rapport de stage en vue d'une soutenance.}

\academicSignSentence{J'autorise l'étudiant à faire le dépôt de son rapport de stage en vue d'une soutenance.}

%================== Résumés ===============================%

\frenchAbstract{
Dans le commerce de détail des télécommunications, le conseiller de vente décide sous contrainte de temps et avec une information partielle, tandis que le responsable de point de vente constate les ruptures de stock plutôt qu'il ne les anticipe. Les outils existants traitent ces deux problèmes séparément, sans les relier ni contextualiser leurs recommandations.

Ce projet conçoit et réalise un moteur agentique qui couple les deux domaines. Il repose sur huit agents intelligents coordonnés par un graphe d'états à mémoire partagée, et sur un principe fondateur : les calculs qui engagent une décision restent déterministes, le modèle de langage formule mais ne calcule pas. La démarche suit le processus de référence CRISP-DM, chaque chapitre du rapport correspondant à une phase du cycle.

Sur un rétro-test de 24 681 prévisions réparties en six plis chronologiques, le modèle appris global réduit de 27,9 pour cent l'erreur du moteur statistique et bat la référence naïve sur 150 des 151 points de vente. Le contrôle de conformité atteint une exactitude parfaite sans aucun faux blocage. L'évaluation en exploitation révèle en revanche un écart de plus d'un point entre la qualité mesurée en laboratoire et la qualité perçue par les utilisateurs, et un taux d'acceptation des propositions inférieur au seuil fixé : le système fonctionne sans encore pleinement convaincre. Ce constat, assumé, oriente les perspectives de travail.
}

\frenchAbstractKeywords{intelligence artificielle agentique, systèmes multi-agents, CRISP-DM, prévision de séries temporelles, optimisation des stocks, génération augmentée par la recherche, validation humaine, commerce de détail}

\englishAbstract{
In telecommunications retail, sales advisors make decisions under time pressure with partial information, while store managers observe stockouts rather than anticipate them. Existing tools address these two problems separately, without connecting them or contextualising their recommendations.

This project designs and implements an agentic engine that couples both domains. It relies on eight intelligent agents coordinated by a shared-state graph, and on one founding principle: computations that commit a decision remain deterministic — the language model phrases, it does not compute. The work follows the CRISP-DM reference process, each chapter of this report corresponding to one phase of the cycle.

On a backtest of 24,681 forecasts across six chronological folds, the global learned model reduces the statistical engine's error by 27.9 percent and beats the naive baseline on 150 of 151 stores. The compliance layer achieves perfect accuracy with no false blocking. Production monitoring, however, reveals a gap of more than one point between laboratory-measured quality and user-perceived quality, and an acceptance rate below the target threshold: the system works without yet fully convincing. This finding, openly stated, shapes the proposed future work.
}

\englishAbstractKeywords{agentic artificial intelligence, multi-agent systems, CRISP-DM, time series forecasting, inventory optimisation, retrieval-augmented generation, human-in-the-loop, retail}

%================== Adresse entreprise ====================%
\setboolean{wantToTypeCompanyAddress}{false}

\companyEmail{contact@ooredoo.tn}
\companyTel{71 000 000}
\companyFax{71 000 001}
\companyAddressAR{تونس}
\companyAddressFR{Tunis, Tunisie}
